% ==========================================================================
% CHAPTER 1 ONLY
% Every technical paragraph is mapped to evidence IDs from
% docs/analysis/evidence_traceability.md and to concrete repository sources.
% Inferences I-001, I-002, I-003, I-024, and I-028 were approved as part of
% the documentation plan; they are phrased as derived framing, not measured
% implementation facts.
% ==========================================================================

\chapter{Introduction}
\label{chap:introduction}

\section{Background and Motivation}
\label{sec:background-motivation}

% Evidence: V-034. Source: RFC 8446.
Secure Internet communication is fundamental to the exchange of personal,
organizational, and governmental information. Transport Layer Security (TLS),
the successor to Secure Sockets Layer (SSL), protects network communication by
providing confidentiality, integrity, and endpoint authentication~\cite{rfc8446}.
In HTTPS, these protections rely on digital certificates that authenticate Web
servers and support the establishment of encrypted connections.

% Evidence: V-034. Source: RFC 5280.
Within a Public Key Infrastructure (PKI), trusted Certificate Authorities (CAs)
issue X.509 certificates that bind identities to public keys. The Internet
profile of X.509 defines certificate fields, extensions, certification paths,
and revocation-related structures~\cite{rfc5280}. A deployed certificate may
therefore describe its subject, issuer, validity period, public-key algorithm
and size, signature algorithm, certificate policies, and Subject Alternative
Names (SANs). Collectively, these attributes provide a basis for evaluating
issuer activity, cryptographic choices, validity practices, domain coverage,
and other characteristics of the certificate ecosystem. As the number of
deployed certificates grows, manual inspection becomes increasingly
impractical, motivating automated collection, structured storage, and analysis.

% Evidence: V-009--V-013, V-034. Source: CT pipeline code/config and RFC 9162.
Certificate Transparency (CT) improves the visibility of certificate issuance
through publicly auditable, append-only logs~\cite{rfc9162}. These logs provide
a valuable source for monitoring issued or observed TLS server certificates;
however, their scale and continuous growth make systematic analysis difficult.
Manual review cannot efficiently track changes in validity, issuance patterns,
or certificate characteristics across extensive and evolving datasets. This
creates a need for analytical approaches that can process large certificate
collections while retaining the contextual information required for meaningful
security assessment.

% Evidence: V-025, V-034. Source: shared-key precompute and committed paper.
Large certificate collections also reveal relationships that are difficult to
identify through isolated certificate lookup. Repeated certificates and reused
RSA public keys have, for example, been observed in large PKI datasets~\cite{nezhadian2025reuse}.
Such relationships may remain hidden when certificates are examined
individually because their significance emerges only after comparing records
across domains, issuers, and time periods. Large-scale analysis is therefore
important for exposing patterns that warrant further investigation, while
ensuring that observations such as key reuse are not treated, by themselves,
as proof of private-key compromise.

% Motivation derived from the approved project scope and problem framing.
The motivation for this project arises from the limitations of manual and
lookup-oriented certificate inspection. Examining one certificate at a time may
answer a narrow operational question, but it does not provide an ecosystem-wide
view of issuer activity, validity practices, cryptographic algorithms, domain
coverage, geographic patterns, or certificate-quality concerns. Researchers,
system administrators, and cybersecurity professionals require a systematic
means of comparing these properties across large datasets so that significant
patterns can be distinguished from isolated observations.

Scalable analytics are necessary because the value of certificate data depends
not only on its availability but also on the ability to organize, summarize,
and interpret it. Raw records and log entries are difficult to assess directly
at scale, whereas clear statistical representations and visual summaries can
make distributions, trends, and unusual relationships easier to investigate.
This need becomes increasingly important as certificate ecosystems continue to
change and as security decisions must be informed by current, comprehensible,
and appropriately contextualized evidence.

These practical challenges motivated the development of the proposed certificate 
analytics platform. The objective is to transform extensive and continuously evolving
SSL/TLS certificate data into organized, meaningful, and actionable information that
supports large-scale security analysis. This naturally leads to the problem addressed
in the following section

\section{Problem Statement}
\label{sec:problem-statement}

% Problem framing approved under I-001; consistent with V-001 and V-009--V-026.
The widespread adoption of HTTPS and the continuing issuance of SSL/TLS
certificates have produced large and continuously changing collections of
certificate data, including records made publicly visible through Certificate
Transparency. Although this information is valuable for assessing the security
posture of Internet services, manually collecting, organizing, and examining it
at scale is impractical. Security-relevant conditions---including expired
certificates, weak cryptographic choices, configuration problems, and public-key
reuse---may also remain difficult to identify when analysis is limited to
individual certificates rather than broader patterns across issuers, domains,
and time periods.

Certificate lookup and monitoring tools support focused inspection, but
individual-record workflows are insufficient for comprehensive, large-scale
security analysis. Researchers, system administrators, and cybersecurity
professionals require a scalable means of organizing and comparing extensive
certificate datasets, transforming raw certificate attributes into
interpretable evidence, and presenting significant distributions and
relationships clearly. Consequently, there is a need for a unified certificate
analytics platform that supports systematic investigation without relying on
manual inspection.

\section{Project Objectives}
\label{sec:objectives}

% Objective framing approved under I-002; each item remains bounded by the
% verified V-series evidence noted in its source comment.
The overall objective of this project is to develop a scalable SSL/TLS
certificate analytics platform that supports the structured collection,
maintenance, analysis, and presentation of large certificate datasets. The
project seeks to make certificate ecosystem data more accessible for systematic
security research, operational monitoring, and informed assessment.

The specific objectives of the project are as follows:

\begin{enumerate}[label=\arabic*.,leftmargin=0.42in]
  \item \textbf{Acquire and organize structured certificate data.} Collect X.509 certificate records from large domain datasets in a form suitable for consistent analysis. % V-004--V-007

  \item \textbf{Maintain an Up-to-Date Certificate Dataset.} Automatically incorporate newly issued and renewed certificates into the analytical dataset to support continuous and up-to-date assessment. % V-009--V-014; freshness is an objective, not a measured outcome

  \item \textbf{Support Multi-Scope Certificate Analytics.} Enable comparative analysis of SSL/TLS certificate ecosystems at both global and country-specific scopes. % V-002--V-003, V-014--V-026

  \item \textbf{Provide Near Real-Time Analytics.} Ensure that newly collected certificate data is reflected in the analytical results with minimal delay, allowing users to monitor continuously evolving certificate ecosystems. % V-004--V-014; timeliness is an objective and is not presented as a measured benchmark

  \item \textbf{Analyze significant certificate characteristics.} Examine validity periods, issuers, cryptographic algorithms, key properties, SANs, and geographic distributions at global and selected country scopes. % V-003, V-014--V-017, V-025--V-026

  \item \textbf{Identify security-relevant patterns.} Surface expiration, weak cryptographic choices, public-key reuse, certificate-quality concerns, and application-defined risk indicators for further investigation. % V-016--V-017, V-025--V-026

  \item \textbf{Support comparative assessment of Certificate Authorities.} Evaluate CA activity and provide comparative rankings using explicitly defined certificate-based metrics. % V-014--V-015; project-defined metrics are not externally calibrated

  \item \textbf{Facilitate interactive certificate analysis.} Present interactive dashboards, statistical summaries, visualizations, search, filtering, and detailed certificate information to facilitate efficient exploration and security analysis. % V-018--V-024
\end{enumerate}

\section{Project Contributions}
\label{sec:contributions}

% Inference: I-003 approved. Verified basis: V-001, V-004--V-026.
The proposed project makes several technical and practical contributions to large-scale
SSL/TLS certificate collection, maintenance, analysis, and visualization. The
principal contributions of the developed platform are summarized below.

\begin{enumerate}[label=\arabic*.,leftmargin=0.42in]
  \item \textbf{Intelligent certificate crawling framework.} The project develops automated crawlers that retrieve and structure SSL/TLS certificates for domains obtained from Certificate Transparency (CT) logs and other supported domain sources. The framework combines concurrent domain-based and CT-driven crawling with a separate experimental IP-based crawler, while persistent task states and stale-task recovery mechanisms provide operational recovery support with limited manual intervention.

  \item \textbf{Continuous certificate dataset maintenance.} A Certificate Transparency driven maintenance process identifies certificates associated with newly observed and renewed domains, incorporates successfully acquired records into the analytical dataset, and supports its continued expansion over time.

  \item \textbf{Timely analytical updates.} Certificate acquisition and dataset maintenance are integrated with analytical processing so that refreshed results can be made available after data-update cycles without requiring a separate manual analysis workflow.

  \item \textbf{Global and country-level analytical scoping.} The platform supports consistent exploration of the complete certificate dataset and configured country-level domain scopes, enabling comparative analysis across multiple geographic scopes within a unified framework.

  \item \textbf{Comprehensive certificate analytics.} The analytical capabilities cover validity periods, issuers, cryptographic algorithms, key properties, validation levels, Subject Alternative Names, and domain-based geographic distributions, providing multiple perspectives on SSL/TLS certificate ecosystems.

  \item \textbf{Security-oriented certificate assessment.} The system surfaces expired certificates, weak cryptographic choices, public-key reuse, certificate-quality findings, and application-defined risk indicators to support focused security investigation.

  \item \textbf{Comparative Certificate Authority evaluation.} The platform evaluates Certificate Authorities using explicitly defined, certificate-derived metrics and presents comparative rankings and distribution patterns.

  \item \textbf{Interactive analysis and visualization.} The web-based interface combines statistical summaries, charts, tables, search, filtering, and detailed certificate views to make large-scale certificate information accessible for systematic exploration.

  \item \textbf{Unified certificate intelligence workflow.} By integrating acquisition, dataset maintenance, analytical processing, and interactive presentation, the project converts raw certificate records into structured and interpretable information for cybersecurity research, operational monitoring, and informed assessment.
\end{enumerate}

\section{Scope and Boundaries}
\label{sec:scope-boundaries}

% Evidence: V-001, V-008, V-014--V-030.
This section defines the functional coverage of the developed platform and the
limits within which its results should be interpreted. It distinguishes the
capabilities addressed by the project from activities and assurances that fall
outside its intended purpose.

\subsection{Project Scope}

The project covers the following principal areas:

\begin{enumerate}[label=\arabic*.,leftmargin=0.42in]
  \item \textbf{Certificate acquisition and dataset maintenance.} Automated crawlers collect SSL/TLS certificates for domains obtained from Certificate Transparency logs and other supported domain sources. Newly observed and renewed certificates extend and update the analytical dataset.

  \item \textbf{Certificate and security analytics.} The platform examines validity, issuers, cryptographic properties, validation levels, Subject Alternative Names, and domain-based geographic distributions. It also identifies expiration, weak cryptographic choices, public-key reuse, certificate-quality findings, and project-defined risk indicators.

  \item \textbf{Certificate Authority analysis.} Certificate Authorities are compared using certificate-derived distributions, quality indicators, and project-defined ranking metrics.

  \item \textbf{Global and country-level analysis.} Certificate data can be examined across the complete dataset and configured country-level domain scopes within a consistent analytical framework.

  \item \textbf{Interactive analysis and visualization.} Dashboards, charts, tables, trends, search, filtering, and detailed certificate views support systematic exploration of analytical results.

  \item \textbf{Analytical refresh.} Analytical results can be refreshed after dataset updates so that subsequent analysis reflects the latest completed acquisition cycle.
\end{enumerate}

\subsection{Project Boundaries}

The project is subject to the following boundaries:

\begin{enumerate}[label=\arabic*.,leftmargin=0.42in]
  \item The dataset does not represent every SSL/TLS certificate deployed on the Internet. Results are limited to certificates acquired from the supported sources.

  \item Security findings are indicators for further investigation, not proof of compromise or confirmed vulnerability. Penetration testing and exploitation of target systems are outside the project scope.

  \item Certificate Authority rankings and risk indicators use project-defined metrics. They are comparative measures rather than official or externally validated security ratings.

  \item Country-level analysis is based on configured domain scopes and does not establish server location, subject jurisdiction, or Certificate Authority location.

  \item Analytical refresh depends on completion of the relevant collection and update cycle. Instantaneous or strictly real-time updates are not claimed.

  \item IP-based certificate acquisition is experimental and remains separate from the primary domain-based analytical dataset.

  \item The platform is intended for research and analytical use; it does not replace production-grade certificate management, threat-intelligence, or formal security-audit services.
\end{enumerate}

Table~\ref{tab:implemented-scope} summarizes the major areas covered by the
project and the principal boundary associated with each area.

\begin{table}[htbp]
  \centering
  \caption{Summary of project scope and principal boundaries}
  \label{tab:implemented-scope}
  \begin{tabularx}{\textwidth}{@{}>{\RaggedRight\arraybackslash\hyphenpenalty=10000\exhyphenpenalty=10000}p{0.19\textwidth}>{\RaggedRight\arraybackslash\hyphenpenalty=10000\exhyphenpenalty=10000}X>{\RaggedRight\arraybackslash\hyphenpenalty=10000\exhyphenpenalty=10000}p{0.31\textwidth}@{}}
    \toprule
    \textbf{Area} & \textbf{Included scope} & \textbf{Principal boundary} \\
    \midrule
    Acquisition & Automated domain-based and CT-guided certificate collection & Coverage is limited to supported sources \\
    Dataset maintenance & Incorporation of newly observed and renewed certificates & Updates depend on completed acquisition cycles \\
    Analytics & Certificate attributes, security indicators, CA comparison, and trends & Findings and metrics require bounded interpretation \\
    Geographic scope & Global and configured country-level domain analysis & Does not establish physical geolocation \\
    Presentation & Interactive summaries, visualizations, search, filtering, and detailed views & Intended for research and analytical use \\
    IP exploration & Experimental IP-based certificate acquisition & Separate from the primary analytical dataset \\
    \bottomrule
  \end{tabularx}
\end{table}

\section{Report Organization}
\label{sec:report-organization}

% Inference: I-024 and I-028 approved as documentation structure.
Following this introduction, Chapter 2 establishes the technical background and compares the three related-work papers committed with the project. Chapter 3 presents repository-derived requirements, constraints, technology choices, and module decomposition. Chapter 4 describes the component, data, scope, indexing, cache, and request-flow design. Chapter 5 examines domain acquisition, dataset preparation, and the experimental IP crawler. Chapter 6 explains the CT renewal and discovery workflow. Chapter 7 documents materialized and live analytics, the Django API, and the Next.js dashboard. Chapter 8 records verification evidence, security observations, implementation limitations, and threats to validity without introducing unsupported benchmark claims. Chapter 9 concludes the report and proposes future work derived from verified repository gaps.
